\documentclass[11pt,letterpaper]{article}
\usepackage[T1]{fontenc}
\usepackage{lmodern,amsmath,amssymb,amsthm,microtype}
\usepackage[margin=1in]{geometry}
\usepackage[hidelinks]{hyperref}
\newtheorem{theorem}{Theorem}
\title{At least four isolates in the height-six Frankl model}
\author{Aafi Khan Malek\\ \small Independent researcher, Ahmedabad, India\\
\small\href{mailto:aafimalek2032@gmail.com}{aafimalek2032@gmail.com}}
\date{28 November 2026}
\begin{document}
\maketitle
A weak family here is finite, nonempty, union-closed, has no empty
member, has no strict inclusion chain of seven members, and has
each point in at most half of its members. Consider a weak family
of minimum ground cardinality. Its proper forcing graph has
an edge \(a\to b\) for distinct points if every member containing
\(a\) also contains \(b\). The earlier
\href{https://github.com/Aafimalek/novel-discoveries/blob/921b03c9fea735147d66af395e1732f01751c62d/frankl-height-six/verification/at-least-three-isolates/paper/few-isolates.pdf}
{few-isolates result} excludes zero, one and two isolates.

\begin{theorem}
The forcing graph of a ground-minimal weak family of height at most six
has at least four isolated points.
\end{theorem}

\section{The pair design at three isolates}
Suppose there are \(k\ge2\) star centers. We use the
\href{https://github.com/Aafimalek/novel-discoveries/tree/043d266655932e3badfe48220085216e90065dd6/frankl-height-six}
{verified structural lemmas}: the stars are unary, the omitted-star
blocks form a pair design on the centers, and for distinct centers
\(p,q\) with common block \(B\), the two-leaf avoidance fiber
\(F_{p,q}\) satisfies
\[
  3k+2+W(p,q)\le |F_{p,q}|\le H(c+3),
  \qquad H(u)=1+u+\binom u2.
\]
Here \(W(p,q)\) counts blocks avoiding both centers, and \(c\) is
the number of distinct occurrence signatures of centers outside
\(B\) in the compressed fiber. The block-link lemma makes
signatures constant on the groups formed by the other blocks
through any center in \(B\).

Write \(N\) for the number of blocks, \(g_v\) for the incidence
degree of center \(v\), \(g=\min_vg_v\), and \(d=k-1-g\).
Choose \(p\) of degree \(g\) and any \(q\ne p\) in a block
through \(p\). The groups through \(p\) give \(c\le g-1\);
\(W(p,q)=N-g-g_q+1\ge N-g-k+2\), so
\[
 N+2k+4-g\le |F_{p,q}|\le H(g+2).
\]
The incidence and Cauchy--Schwarz bound from the earlier proof,
\(N\ge kg^2/(k-1+g)\), gives after substitution and multiplication
by the positive denominator
\[
 (d-4)g^2+(5d+8)g+4d(d+1)\le0 .
\]
Since \(g\ge1\), this forces \(d\le3\). A size-\(r\) block
contributes \(r(r-2)\) to the sum of incidence deficits and replaces
\(\binom r2\) pairs by one block. Since
\(\binom r2-1\le 2r(r-2)/3\), summing gives
\(N\ge\binom k2-2kd/3\). Inserting this lower bound into the fiber
comparison gives at most
\[
   (1-d/3)k+(d^2-5d-6)/2
\]
units of slack. It equals \(-6\) at \(d=3\), so \(d\le2\).

If \(d=1\) or 2, a group through \(q\) cannot meet two distinct
groups through \(p\). If it did, their signatures would agree and
the upper fiber bound would drop by \(H(g+2)-H(g+1)=g+2\);
this exceeds the slack \(2k/3-5\) for \(d=1\), or \(k/3-6\)
for \(d=2\). Thus each \(q\)-group is contained in a \(p\)-group.
If a block through \(q\) other than \(B\) had size at least three,
two of its outside points would lie in one \(p\)-group, violating
uniqueness of the block containing their pair. This applies to
every \(q\ne p\) (with its unique block through \(p\)), so every
nonpair block contains \(p\). With \(d=1\) there is exactly one
triple and \(N=\binom k2-2\); the lower fiber bound is then
\(\binom k2+k+4>H(k)\), its upper bound. With \(d=2\),
there are two triples or one four-block, giving
\(N\ge\binom k2-5\), so the lower bound is
\(\binom k2+k+2>H(k)\), already larger than the actual upper
bound \(H(k-1)\). Only \(d=0\) remains. All blocks are pairs.

\section{Relative co-singletons and the contradiction}
Now \(N=\binom k2\), and every leaf-pair fiber has size at
least \(N+k+5\). The compression retains at most the other
\(k-2\) stars and the three isolates. Its support has all
\(n=k+1\) coordinates: if one were always absent, the upper
bound would be just \(H(k)=N+k+1\). By union closure the full
support is present. The fiber has height at most three, and
\(N+k+5\le H(k+1)\) also implies \(k\ge3\).

Here is a counting consequence of the small gap to \(H(n)\).
On the full support take complements of the compressed fiber.
They form an intersection-closed family \(\mathcal D\) containing
\(\varnothing\), also of height at most three. Let \(a\) be the
number of its minimal nonempty members, called atoms. They are
disjoint. Every other nonempty member is maximal and contains an
atom. A maximal nonatom containing two or more atoms consumes
a pair of atoms that cannot occur in another maximal nonatom.
A maximal nonatom containing exactly one atom must contain a
point outside the union of the atoms; no point of this kind can
belong to two maximal nonatoms, since their nonempty intersection
would have to be an atom. There are at most \(n-a\) outside
points. Thus
\[
  |\mathcal D|\le1+a+\binom a2+(n-a)=1+n+\binom a2.
\]
If \(a\le n-1\), the right side is at most \(N+k+2\), a
contradiction. Hence all \(n\) singletons belong to
\(\mathcal D\). In the original family, for each pair of stars,
the maximum avoiding their leaves and every relative co-singleton
of that maximum consequently occur.

Write \(U\) for the full ground set and \(S_r\) for the unary
star with center \(r\), including its leaf. For a center \(r\) and
isolate \(x\), choose two other distinct centers \(a,b\).
The relative co-singletons just proved give the two members with
omitted coordinates \(S_r\cup S_a\cup\{x\}\) and
\(S_r\cup S_b\cup\{x\}\). Their union is the member with
complement \(S_r\cup\{x\}\).

Fix two leaf anchors at distinct centers \(p,q\). These give
\(3(k-2)\) new both-containing members, using \(r\ne p,q\)
and any isolate \(x\). They are pairwise distinct; they omit
three actual points, unlike the canonical top and single-point
avoidance maxima. They omit a leaf, unlike the source-witness
maxima, and omit only one star, unlike the pair maxima. Therefore
the same verified noncanonical-witness counting inequality adds
\(3(k-2)\) to the lower bound \(N+k+5\). But the upper bound
\(H(k+1)\) leaves only \(k-3\) available. This contradiction
excludes \(k\ge2\).

\section{The one-center structure}
If there is no center then the whole ground consists of the three
isolates. No singleton is possible in a ground-minimal weak
family, so each member would contain at least two of three
points, contradicting the half-frequency bounds.

With one center \(p\), every member containing it contains its
entire star. The existing containing-fiber lemma supplies a point
\(x\) occurring in strictly fewer than half these members. It must
be an isolate. With \(y,z\) the other isolates, the full isolate
trace and all two-element traces occur by the co-singleton lemma.
There are already three \(x\)-containing traces. Having fewer
than half forces all four \(x\)-free traces and excludes the
optional trace \(\{x\}\). Thus exactly seven members contain
\(p\), including the pure star. Union with the pure star
excludes trace \(\{x\}\) everywhere.

Write \(\mathcal B_0\) for the \(b\) members avoiding \(p\)
and \(\ell\) for the number of leaves. The frequencies of \(y,z\)
in \(\mathcal B_0\) are each at most \((b-1)/2\), because each
occurs in four of the seven containing members. If \(a_0\) and
\(a_2\) are the counts in \(\mathcal B_0\) omitting both or
containing both of \(y,z\), respectively, inclusion--exclusion gives
\[
    a_0=b-f_{\mathcal B_0}(y)-f_{\mathcal B_0}(z)+a_2
       \ge a_2+1.
\]
The full avoiding maximum \(U\setminus\{p\}\) and all
\(U\setminus\{p,\lambda\}\) for leaves \(\lambda\) belong to
\(\mathcal B_0\): \(p\) has no proper predecessor, and is the
unique proper predecessor of each leaf. Each of these \(\ell+1\)
distinct sets contains \(y,z\), so \(a_0\ge\ell+2\).

Since the trace \(\{x\}\) is forbidden, the members counted by
\(a_0\) have no isolates at all. They form a union-closed subfamily
on the leaves. It has height at most two: any three members in
a strict chain can be followed by the four distinct containing
members whose isolate traces are
\(\varnothing,\{y\},\{y,z\},\{x,y,z\}\). The first of these four
is strictly above every member on the leaves alone, since it
contains \(p\). The height-two bound applied on at most
\(\ell\) coordinates gives \(a_0\le1+\ell\), a contradiction.
The height-two bound also follows directly: every member other
than the maximum misses a point not missed by any other such
member, since the union of two different nonmaximum members
must be the maximum.

This excludes three isolates in all cases. Along with the earlier
exclusions of zero, one and two isolates, this proves the theorem.

\paragraph{Scope and contributions.}
The argument is written mathematics using previously Lean-checked
structural lemmas. Lean checks the indicated arithmetic sublemmas separately; the new composition is not yet checked end to end. The case of four or more isolates remains open.
Aafi Khan Malek selected and directed the research. ChatGPT assisted
with proof development and writing. Codex was not used.
This note has not been peer reviewed.
\end{document}
